% ═══════════════════════════════════════════════════════════════════════════════
% Sections 1–3 for PRICAI submission
% Springer LNCS format  |  double-blind
% ═══════════════════════════════════════════════════════════════════════════════

% ─── Preamble (merge into your main .tex file) ────────────────────────────────
\documentclass[runningheads]{llncs}
\usepackage{amsmath,amssymb}
\usepackage{graphicx}
\usepackage{booktabs}
\usepackage{xcolor}
\usepackage{hyperref}
\usepackage[T1]{fontenc}

\begin{document}

\title{Learning Multi-Scale Market Representations\\
       for Short-Horizon Cryptocurrency Trading}

\author{Anonymised for review}
\institute{Anonymised for review}
\maketitle

% ─────────────────────────────────────────────────────────────────────────────
\begin{abstract}
We propose \textsc{HybridSSM-Trader}, a multi-scale temporal backbone for short-horizon
cryptocurrency price prediction that combines Temporal Convolutional Networks
(TCN), parallel Mamba state-space branches, and multi-head self-attention through
a learned three-way gated fusion.  Crucially, \textsc{HybridSSM-Trader} is trained with
a multi-task objective that jointly optimises a price regression head and an
action classification head (sell/hold/buy), and we show experimentally that this
joint objective is \emph{not} a minor add-on: removing it causes approximately
40\% of random seeds to converge to an \emph{anti-predictive} signal, collapsing
directional accuracy below 15\%.  On a frozen Bitcoin daily test set spanning
January~2025 to March~2026, the full hybrid achieves a test-RMSE standard
deviation of $\pm 4$ across five seeds, compared with $\pm 14$--$15$ for
single-branch TCN and Mamba baselines and $\pm 469$ for an attention-only model,
demonstrating substantially more stable generalisation.  Component ablations
confirm that each of the three temporal branches and the gated fusion
contribute to this stability.
\end{abstract}

% ─────────────────────────────────────────────────────────────────────────────
\section{Introduction}
\label{sec:intro}

Short-horizon cryptocurrency forecasting sits at the intersection of two
well-established difficulties: financial markets exhibit near-efficient pricing
dynamics~\cite{fama1970efficient}, and high-frequency price series are
non-stationary, heavy-tailed, and subject to abrupt regime changes.  Deep neural
networks have achieved strong results on structured sequence tasks, yet their
application to financial prediction is constrained by a fundamental mismatch
between the optimisation objective and the downstream decision task.  Mean
squared error---the default regression loss---penalises magnitude errors uniformly
and is entirely indifferent to whether the \emph{direction} of a predicted price
move is correct.  For a trading agent, a small directional error that induces a
wrong buy/sell decision is far more costly than a large magnitude error in a
correctly-signed prediction.

A second challenge is representational.  Financial prices carry structure at
multiple temporal scales simultaneously: high-frequency patterns (momentum,
mean-reversion over a few days), medium-horizon trends (trend persistence over
weeks), and global context (macro regime, crowd sentiment reflected across the
full observation window).  Existing work largely addresses these scales in
isolation.  TCN-based models~\cite{bai2018empirical} excel at capturing local
dilated patterns but lack a mechanism for long-range state propagation.
Transformer variants~\cite{zhou2021informer,wu2021autoformer} attend globally but
are known to suffer from instability on small, noisy financial datasets.
Mamba~\cite{gu2023mamba}, a selective state-space model, offers efficient linear
recurrence well-suited to long-horizon state tracking but cannot directly
express local dilated features.

We propose \textsc{HybridSSM-Trader}, a backbone that processes each observation window
through all three families in \emph{parallel} and fuses their outputs through a
learned, position-wise three-way softmax gate.  A volatility-gated exponential
smoother is applied as a post-processing step to attenuate noise adaptively under
different market regimes.  Importantly, rather than relying on a single price
regression objective, we train the model jointly to predict both the next-period
price and a ternary trading action (sell/hold/buy) derived from the predicted
return magnitude relative to a transaction-cost-adjusted threshold.  Our
experiments reveal that this multi-task objective is essential for learning
stable directional representations: without it, the backbone has no gradient
signal that distinguishes a correct direction from an incorrect one, and
converges to the regression optimum regardless of signal polarity.

\paragraph{Contributions.}
\begin{enumerate}
  \item \textbf{Multi-scale backbone with corrected fusion.}
        We introduce \textsc{HybridSSM-Trader}, a hybrid TCN--Mamba--Attention
        architecture with a three-way \emph{softmax} gated feature fusion.
        We identify and correct a systematic flaw present in concurrent
        work~\cite{liu2025alphaseek}: their sigmoid-based gate assigns a
        fixed scalar weight to the global branch regardless of input, which
        prevents the gate from learning when attention context is uninformative.
        Our formulation eliminates this bias; our ablation (``simple fusion'')
        quantifies the resulting performance gap.
  \item \textbf{Multi-task stabilisation.}  We demonstrate through a controlled
        ablation that the joint action-classification objective is necessary for
        reliable directional learning: removing it causes $\sim$40\% of random
        initialisations to converge to anti-predictive signals
        (directional accuracy $< 0.15$, worse than chance).
  \item \textbf{Rigorous frozen-test evaluation.}  We report results across five
        random seeds on a completely held-out BTC test set (Jan 2025--Mar 2026)
        that was not used during any model selection decision, providing an
        unusually clean out-of-sample evaluation for cryptocurrency ML research.
\end{enumerate}

% ─────────────────────────────────────────────────────────────────────────────
\section{Related Work}
\label{sec:related}

\paragraph{Deep learning for financial time series.}
Early applications of recurrent networks to price prediction~\cite{hochreiter1997long}
established that learned sequence models can capture autocorrelation in financial
returns, though reliable generalisation requires careful regularisation.  TCNs,
introduced for sequence modelling by Bai et al.~\cite{bai2018empirical}, use
dilated causal convolutions to capture long-range dependencies without recurrence
and have since been applied to intraday equity and cryptocurrency
prediction~\cite{sezer2020financial}.  Transformer-based architectures have
been applied to multivariate financial time series~\cite{zhou2021informer},
but attention scales quadratically with sequence length and is prone to instability
on datasets with high noise-to-signal ratios typical of cryptocurrency markets.

\paragraph{State-space models.}
Structured state-space models (S4~\cite{gu2022efficiently}) address the
computational limitations of attention on long sequences through a linear
recurrence parameterised by a structured transition matrix.  Mamba~\cite{gu2023mamba}
extends S4 with input-selective state transitions, enabling the model to focus
selectively on informative timesteps.  Recent work has demonstrated promising
results applying Mamba to financial time series~\cite{xu2024mamba}, though these
studies typically treat the Mamba backbone in isolation rather than as one
component of a multi-scale ensemble.

\paragraph{Multi-scale and hybrid architectures.}
TimesNet~\cite{wu2023timesnet} and PatchTST~\cite{nie2022time} have explored
the combination of local patch-level and global sequence-level representations,
motivating the value of explicit multi-scale design.  Closest to our architecture
is the ``Crossformer''~\cite{zhang2023crossformer}, which uses a two-stage
attention over time and feature dimensions; however, it does not include a Mamba
branch or a volatility-adaptive smoothing stage.  Concurrent work on hybrid
SSM-attention models for general time series~\cite{lieber2024jamba} shows that
the two families are complementary, supporting our parallel fusion approach.

\paragraph{Closest prior work.}
Liu et al.~\cite{liu2025alphaseek} independently propose a backbone combining
Mamba SSM, TCN, and multi-head attention for high-frequency cryptocurrency
trading under the name ``AlphaSeek FinRL''.  Despite architectural similarity,
three substantive differences separate their work from ours.
First, their gated fusion uses a binary sigmoid gate
$\mathbf{f} = \sigma(\cdot)\!\cdot\!\mathbf{h}_\text{local}
 + (1{-}\sigma(\cdot))\!\cdot\!\mathbf{h}_\text{long}
 + 0.5\!\cdot\!\mathbf{h}_\text{global}$:
the global (attention) branch receives a \emph{fixed} additive scalar~(0.5)
that is independent of any input signal, preventing the gate from learning
how much global context is actually useful.
\textsc{HybridSSM-Trader} replaces this with a three-way softmax gate
(Eq.~\ref{eq:gate}) in which all branches compete under a probability simplex,
eliminating the free-riding bias.
Second, Liu et al.\ apply sequential BatchNorm then LayerNorm input
preprocessing---a redundant double-normalisation that our ablation confirms
is unnecessary; we use a single LayerNorm (Eq.~\ref{eq:preprocess}).
Third, and most critically, their architecture has no action classification
head and is trained solely on an unspecified regression objective; we show
in \S\ref{sec:results} that this omission leads to training instability in
$\sim$40\% of random seeds.  Finally, their evaluation uses a single random
seed with no held-out test split, whereas we report five-seed statistics on
a frozen test set, providing a substantially more rigorous basis for comparison.

\paragraph{Trading objectives and multi-task learning.}
The disconnect between MSE minimisation and downstream trading performance is
well-documented~\cite{fischer2018deep}.  Several works have proposed direction-
aware loss functions, including the Mean Absolute Directional
Loss~\cite{michalikow2023madl} and its differentiable extension
GMADL~\cite{sakowski2024gmadl}, which penalise magnitude-weighted directional
errors rather than raw forecast deviations.  Multi-task formulations that
combine price regression with action classification have been explored in
equity trading~\cite{ding2015deep}, but the effect of the classification head on
the stability of directional learning has not been directly measured.  Our
ablation study fills this gap.

\paragraph{Cryptocurrency-specific forecasting.}
Bitcoin price prediction has attracted dedicated study due to its high
volatility and 24/7 trading~\cite{ji2019does}.  Regime-aware models and
walk-forward evaluation protocols are increasingly recognised as necessary for
credible evaluation~\cite{atsalakis2019bitcoin}, motivating our use of a
completely frozen test split with explicit train/validation/test temporal
boundaries.

% ─────────────────────────────────────────────────────────────────────────────
\section{Method}
\label{sec:method}

\subsection{Problem Formulation}
\label{sec:formulation}

Let $\mathbf{x}_t \in \mathbb{R}^F$ denote the feature vector at day $t$,
comprising OHLCV prices and six engineered features (close return, log-return,
high--low range, open--close change, volume change, and seven-day rolling
volatility), giving $F=12$.  Given a window of $T=30$ consecutive daily
observations $\mathbf{X} = [\mathbf{x}_{t-T+1}, \ldots, \mathbf{x}_t]
\in \mathbb{R}^{F \times T}$, the model is asked to predict both (i) the
next-day close price $y_{t+1}$ and (ii) a ternary trading action
$a_{t+1} \in \{0\text{(sell)},\, 1\text{(hold)},\, 2\text{(buy)}\}$.

The action label is derived from the one-step return
$r_{t+1} = (y_{t+1} - y_t) / y_t$ relative to a threshold $\tau_{\text{eff}}
= \tau + c_{\text{tx}}$, where $\tau$ is the minimum return threshold and
$c_{\text{tx}}$ is the round-trip transaction cost:
\begin{equation}
  a_{t+1} =
  \begin{cases}
    2 & r_{t+1} >  +\tau_{\text{eff}} \\
    0 & r_{t+1} <  -\tau_{\text{eff}} \\
    1 & \text{otherwise.}
  \end{cases}
  \label{eq:action_label}
\end{equation}

\subsection{Input Preprocessing}
\label{sec:preprocess}

The raw input $\mathbf{X} \in \mathbb{R}^{F \times T}$ is first transposed to
$\mathbb{R}^{T \times F}$ and normalised by a per-feature LayerNorm:
\begin{equation}
  \mathbf{H}_0 = \text{Dropout}\!\left(
    \text{GELU}\!\left(
      \mathbf{W}_p\, \text{LN}(\mathbf{X}^\top)
    \right)
  \right) \in \mathbb{R}^{T \times d},
  \label{eq:preprocess}
\end{equation}
where $\mathbf{W}_p \in \mathbb{R}^{F \times d}$ is a learned projection and
$d = 64$ is the shared hidden dimension.  A single LayerNorm over the feature
axis is used, as it is the standard normalisation for position-wise sequence
models and avoids the redundant double-normalisation common in prior
financial time series models.

\subsection{Parallel Temporal Branches}
\label{sec:branches}

\textsc{HybridSSM-Trader} processes $\mathbf{H}_0$ through three specialised
encoders in parallel, each producing a context tensor of the same shape
$\mathbb{R}^{B \times T \times d}$.

\paragraph{TCN branch (local context).}
A two-layer Temporal Convolutional Network with kernel size $k=3$, exponentially
increasing dilation ($2^0, 2^1$), and weight normalisation captures local
pattern dynamics:
\begin{equation}
  \mathbf{H}_{\text{local}} = \text{TCN}(\mathbf{H}_0).
  \label{eq:tcn}
\end{equation}
Each temporal block applies two dilated causal convolutions with GELU activations
and a skip connection with a $1{\times}1$ projection if the channel dimension changes.

\paragraph{Mamba branch (long-range context).}
The hidden representation is split into $K=4$ equal-width chunks of dimension
$d/K$ and processed by $K$ independent Mamba selective state-space
modules~\cite{gu2023mamba}:
\begin{equation}
  \mathbf{H}_{\text{long}} = \text{LN}\!\left(
    \mathbf{H}_0 +
    \text{Dropout}\!\left(
      \bigoplus_{k=1}^{K} \text{Mamba}_k\!\left(\mathbf{H}_0^{(k)}\right)
    \right)
  \right),
  \label{eq:mamba}
\end{equation}
where $\bigoplus$ denotes concatenation along the feature axis and
$\mathbf{H}_0^{(k)}$ is the $k$-th chunk of a learned pre-projection of
$\mathbf{H}_0$.  Each Mamba module uses state dimension $d_s{=}16$,
convolutional kernel $d_c{=}4$, and expansion factor $e{=}2$.
The multi-branch design allows parallel branches to attend to different
temporal patterns within the state-space parameterisation.

\paragraph{Attention branch (global context).}
A single multi-head self-attention layer with $h=4$ heads captures global
pairwise relationships across all positions:
\begin{equation}
  \mathbf{H}_{\text{global}}, \;\mathbf{A}
    = \text{MHA}\!\left(\bar{\mathbf{H}}_0,\;
                        \bar{\mathbf{H}}_0,\;
                        \bar{\mathbf{H}}_0\right),
  \quad \bar{\mathbf{H}}_0 = \text{LN}(\mathbf{H}_0),
  \label{eq:attention}
\end{equation}
where the normalised representation $\bar{\mathbf{H}}_0$ is computed once and
shared as query, key, and value to avoid redundant computation.
$\mathbf{A} \in \mathbb{R}^{B \times h \times T \times T}$ is retained for
analysis.

\subsection{Gated Feature Fusion}
\label{sec:fusion}

The three context tensors are combined through a position-wise three-way softmax
gate.  Each branch is first projected into a shared hidden space:
\begin{align}
  \mathbf{p}_j &= \mathbf{W}_j^P\, \mathbf{H}_j, \quad j \in
    \{\text{local}, \text{long}, \text{global}\}, \\
  \mathbf{w} &= \text{softmax}\!\left(
    \text{MLP}_{gate}\!\left(
      [\mathbf{p}_{\text{local}} \;\|\; \mathbf{p}_{\text{long}}
       \;\|\; \mathbf{p}_{\text{global}}]
    \right)
  \right) \in \mathbb{R}^{B \times T \times 3},
  \label{eq:gate}
\end{align}
where $\mathbf{W}_j^P \in \mathbb{R}^{d \times d}$ and $\text{MLP}_{gate}$
maps $3d \to d \to 3$ with GELU activation.  The softmax normalisation ensures
the three weights sum to one at every position, so each branch must earn its
contribution and no branch receives a fixed additive bias.  The fused
representation is:
\begin{equation}
  \mathbf{H}_{\text{fused}} =
    w_1 \mathbf{p}_{\text{local}}
    + w_2 \mathbf{p}_{\text{long}}
    + w_3 \mathbf{p}_{\text{global}},
  \label{eq:fused}
\end{equation}
followed by a linear output projection, a residual addition from the input
$\mathbf{H}_0$, and a LayerNorm.  A two-layer position-wise FFN ($d \to 128
\to d$) with GELU is then applied to the fused sequence.

\subsection{Adaptive Signal Smoother}
\label{sec:smoother}

Raw hidden states can be noisy under high-volatility market conditions.
We apply a learned exponential moving average whose decay parameter is
conditioned on the current-window close-price volatility:
\begin{align}
  \hat{\sigma}_v &= \text{std}\!\left(
    \frac{\Delta p_{\text{close}}}{p_{\text{close}}}
  \right)_{\,t-T+1}^{\,t}, \label{eq:vol} \\
  \alpha &= \alpha_{\min} + (\alpha_{\max} - \alpha_{\min})
            \cdot \sigma\!\left(\mathbf{w}_v^\top \hat{\sigma}_v\right),
            \label{eq:alpha} \\
  \tilde{\mathbf{h}}_t &= \alpha \, \mathbf{h}_t
                         + (1 - \alpha)\, \tilde{\mathbf{h}}_{t-1},
  \label{eq:ema}
\end{align}
where $\mathbf{w}_v \in \mathbb{R}^d$ is a learned gating vector,
$\alpha_{\min}{=}0.15$ and $\alpha_{\max}{=}0.85$.  High volatility yields
$\alpha \to \alpha_{\max}$, giving more weight to the current step;
low volatility yields heavier smoothing.  The smoother is implemented without
a Python loop via vectorised causal decay weights, enabling efficient
GPU execution.

\subsection{Multi-Task Signal Heads}
\label{sec:heads}

The final hidden state of the last timestep, $\mathbf{h}_{-1} \in \mathbb{R}^d$,
is passed to two independent heads after a LayerNorm:
\begin{align}
  \hat{y} &= \mathbf{W}_2^R \;\text{GELU}(\mathbf{W}_1^R \mathbf{h}_{-1})
             \in \mathbb{R},
             \label{eq:reghead} \\
  \hat{\mathbf{a}} &= \mathbf{W}_2^A \;\text{GELU}(\mathbf{W}_1^A \mathbf{h}_{-1})
                    \in \mathbb{R}^3,
             \label{eq:actionhead}
\end{align}
where $\mathbf{W}_1^R, \mathbf{W}_1^A \in \mathbb{R}^{d/2 \times d}$ and
$\mathbf{W}_2^R \in \mathbb{R}^{1 \times d/2}$,
$\mathbf{W}_2^A \in \mathbb{R}^{3 \times d/2}$.

\subsection{Training Objective}
\label{sec:loss}

The total loss is a weighted combination of a regression term and a
classification term:
\begin{equation}
  \mathcal{L} = \lambda_r \mathcal{L}_{\text{reg}} +
                \lambda_a \mathcal{L}_{\text{CE}},
  \label{eq:loss}
\end{equation}
where $\mathcal{L}_{\text{reg}} = \sqrt{\text{MSE}(\hat{y}, y)}$ (RMSE),
$\mathcal{L}_{\text{CE}}$ is the cross-entropy loss over the three action classes,
$\lambda_r = 1.0$, and $\lambda_a = 0.05$.  The regression target is the
raw next-day close price (no normalisation applied at the dataset level), and
both $\hat{y}$ and $y$ are in the same price space.

\paragraph{Why the action head matters.}
The regression loss penalises magnitude error but carries no gradient signal about
prediction \emph{direction}: a prediction of $\hat{y} = y - \epsilon$ incurs
almost no loss when $\epsilon$ is small, even if the sign of the move is wrong.
The action cross-entropy fills this gap: correctly classifying a buy day as buy
(rather than sell) contributes a gradient that shapes the representation toward
correct directional encoding.  We confirm this in \S\ref{sec:results}: removing
$\mathcal{L}_{\text{CE}}$ entirely causes approximately 40\% of random seeds to
converge to a model with directional accuracy below 15\%---worse than random
assignment over three classes---indicating that the regression objective alone
permits anti-correlated signal solutions.  The action weight $\lambda_a = 0.05$
is chosen to be small enough that the primary optimisation target remains price
accuracy while the classification term acts as a directional regulariser.

\paragraph{Model summary.}  Table~\ref{tab:arch} summarises the architecture
hyperparameters used in all experiments reported in this paper.

\begin{table}[h]
\centering
\caption{HybridSSM-Trader architecture hyperparameters.}
\label{tab:arch}
\begin{tabular}{lc}
\toprule
\textbf{Component} & \textbf{Value} \\
\midrule
Input window $T$ & 30 \\
Input features $F$ & 12 \\
Hidden dimension $d$ & 64 \\
TCN layers / kernel / dilation & 2 / 3 / $2^i$ \\
Mamba branches $K$ & 4 \\
Mamba state dim $d_s$ & 16 \\
Attention heads $h$ & 4 \\
FFN hidden dim & 128 \\
Smoother $[\alpha_{\min}, \alpha_{\max}]$ & $[0.15, 0.85]$ \\
Action classes & 3 \\
\midrule
Optimiser & AdamW \\
Learning rate & $5 \times 10^{-4}$ \\
LR schedule & CosineAnnealingLR ($T_{\\max}{=}100$) \\
Weight decay & $10^{-4}$ \\
Batch size & 32 \\
Early stopping patience & 30 epochs (val RMSE) \\
\bottomrule
\end{tabular}
\end{table}

% ─────────────────────────────────────────────────────────────────────────────
% Placeholder bibliography entries — replace with your .bib file
% ─────────────────────────────────────────────────────────────────────────────
\bibliographystyle{splncs04}
\begin{thebibliography}{99}

\bibitem{fama1970efficient}
Fama, E.F.: Efficient capital markets: A review of theory and empirical work.
J. Finance \textbf{25}(2), 383--417 (1970)

\bibitem{bai2018empirical}
Bai, S., Kolter, J.Z., Koltun, V.: An empirical evaluation of generic convolutional
and recurrent networks for sequence modeling. arXiv:1803.01271 (2018)

\bibitem{hochreiter1997long}
Hochreiter, S., Schmidhuber, J.: Long short-term memory.
Neural Comput. \textbf{9}(8), 1735--1780 (1997)

\bibitem{zhou2021informer}
Zhou, H., et al.: Informer: Beyond efficient transformer for long sequence
time-series forecasting. In: AAAI, pp. 11106--11115 (2021)

\bibitem{wu2021autoformer}
Wu, H., et al.: Autoformer: Decomposition transformers with auto-correlation
for long-term series forecasting. In: NeurIPS (2021)

\bibitem{gu2023mamba}
Gu, A., Dao, T.: Mamba: Linear-time sequence modeling with selective state spaces.
arXiv:2312.00752 (2023)

\bibitem{gu2022efficiently}
Gu, A., et al.: Efficiently modeling long sequences with structured state spaces.
In: ICLR (2022)

\bibitem{xu2024mamba}
Xu, M., et al.: MambaMixer: Efficient selective state space models with dual
token and channel selection. arXiv:2403.19888 (2024)

\bibitem{wu2023timesnet}
Wu, H., et al.: TimesNet: Temporal 2D-variation modeling for general time series
analysis. In: ICLR (2023)

\bibitem{nie2022time}
Nie, Y., et al.: A time series is worth 64 words: Long-term forecasting with
transformers. In: ICLR (2023)

\bibitem{zhang2023crossformer}
Zhang, Y., Yan, J.: Crossformer: Transformer utilizing cross-dimension dependency
for multivariate time series forecasting. In: ICLR (2023)

\bibitem{lieber2024jamba}
Lieber, O., et al.: Jamba: A hybrid transformer-mamba language model.
arXiv:2403.19887 (2024)

\bibitem{fischer2018deep}
Fischer, T., Krauss, C.: Deep learning with long short-term memory networks
for financial market predictions. Eur. J. Oper. Res. \textbf{270}(2), 654--669
(2018)

\bibitem{michalikow2023madl}
Michańków, J., et al.: Mean absolute directional loss as a new loss function
for machine learning problems in algorithmic investment strategies.
arXiv:2309.10546 (2023)

\bibitem{sakowski2024gmadl}
Sakowski, P., et al.: Generalized mean absolute directional loss as a solution
to overfitting and high transaction costs in machine learning models used in
high-frequency algorithmic investment strategies. arXiv:2412.18405 (2024)

\bibitem{ding2015deep}
Ding, X., et al.: Deep learning for event-driven stock prediction.
In: IJCAI, pp. 2327--2333 (2015)

\bibitem{sezer2020financial}
Sezer, O.B., et al.: Financial time series forecasting with deep learning:
A systematic literature review. Appl. Soft Comput. \textbf{90}, 106181 (2020)

\bibitem{ji2019does}
Ji, Q., et al.: Does Bitcoin contribute to hedge against inflation?
J. Int. Financ. Mark. \textbf{62}, 1--9 (2019)

\bibitem{atsalakis2019bitcoin}
Atsalakis, G.S., et al.: Bitcoin price forecasting with neuro-fuzzy techniques.
Eur. J. Oper. Res. \textbf{276}(2), 770--780 (2019)


\bibitem{liu2025alphaseek}
Liu, J., Ma, J., Jiang, Z.: AlphaSeek FinRL: A hybrid deep learning architecture
for high-frequency cryptocurrency trading. In: IEEE 11th International Conference
on Intelligent Data and Security (IDS), pp.~54--56 (2025)

\end{thebibliography}

\end{document}
